
{\Large\bfseries\filcenter
Data Management and Sharing Plan (DMSP)
\\\rule{\textwidth}{0.4pt}}

\setcounter{section}{0}
\section{Expected Data}

During the course of the project, the following data will be produced.
\begin{enumerate}
    \item Various robotic datasets, benchmarks, examples, and documentations/reports on their purpose details.
    \item Source code, tools, tutorials, demos on robotic models and their security and properties.
    \item Educational material, lectures notes, lab experiments and activities.
    \item Research papers, pre-prints, technical reports.
\end{enumerate}

\section{Data Format}

Most of the documentation, research, and educational data are expected to be in textual format, \eg code, reports, notes, lectures, and scripts. 
Public dataset, benchmarks, and models will be available on GitHub and the local server. 
Lecture materials will be in pdf format.

\section{Access to Data and Data Sharing Practices and Policies}

A website for the project will be set up that will allow data sharing. Research and educational materials, including code, datasets, demos, lectures, and labs will be shared with the research community. 
We commit to releasing the project's defense tools, benchmarks, datasets, and source code as open source under permissive licenses to support reproducibility and community reuse.
Offensive attack tooling follows a staged, controlled-access release governed by data-use agreements and responsible-disclosure review, keeping defensive artifacts openly available while dual-use attack code is gated to prevent misuse.
The robotic security course will be posted to the CLARK course library.
%Spotlights on each participant will also be included. Many participants will likely put any coding or scripts, and eventually their certificate on their Github or LinkedIn profile to showcase their work. Aggregate results, such as demographic information, can be shared.

%All data will be stored in a Microsoft OneDrive folder, making access control easier to manage.

\section{Policies for Re-Use, Re-Distribution}

Any of the public information on the program website can be re-used and re-distributed, as long as appropriate citations are made. Access to raw data could be provided upon request as long as no private information is included.

\section{Archiving of Data}

All data produced will be archived after the end of the program. The data will be kept according to Loyola University Chicago archival policies and NSF guidelines and securely deleted after that archival period, or after five years, whichever is earlier. %Aggregated data regarding the participants' information will be kept for longer for possible publication or submission for a workshop continuation proposal. This does not contain any private information.



% Proposals must include a supplementary document of no more than two pages labeled ``Data Management Plan". This
% supplementary document should describe how the proposal will conform to NSF policy on the
% dissemination and sharing of research results (see AAG Chapter VI.D.4)
